\documentclass[11pt]{article}

\usepackage[margin=1in]{geometry}
\usepackage[T1]{fontenc}
\usepackage[utf8]{inputenc}
\usepackage{lmodern}
\usepackage{microtype}
\usepackage{amsmath,amssymb,amsthm,mathtools}
\usepackage{booktabs}
\usepackage{enumitem}
\usepackage{tabularx}
\usepackage[hidelinks]{hyperref}
\usepackage[nameinlink,capitalise]{cleveref}

\newcolumntype{L}[1]{>{\raggedright\arraybackslash}p{#1}}
\newcolumntype{Y}{>{\raggedright\arraybackslash}X}

\newtheorem{theorem}{Theorem}[section]
\newtheorem{proposition}[theorem]{Proposition}
\newtheorem{lemma}[theorem]{Lemma}
\newtheorem{corollary}[theorem]{Corollary}
\theoremstyle{definition}
\newtheorem{definition}[theorem]{Definition}
\newtheorem{assumption}[theorem]{Assumption}
\newtheorem{example}[theorem]{Example}
\theoremstyle{remark}
\newtheorem{remark}[theorem]{Remark}

\newcommand{\cA}{\mathcal A}
\newcommand{\cD}{\mathcal D}
\newcommand{\cH}{\mathcal H}
\newcommand{\cI}{\mathcal I}
\newcommand{\cK}{\mathcal K}
\newcommand{\cM}{\mathcal M}
\newcommand{\cP}{\mathcal P}
\newcommand{\cQ}{\mathcal Q}
\newcommand{\cR}{\mathcal R}
\newcommand{\cX}{\mathcal X}
\newcommand{\Id}{\operatorname{Id}}
\newcommand{\dist}{\operatorname{dist}}
\newcommand{\spec}{\operatorname{spec}}
\newcommand{\MTT}{Modal Triplet Theory}

\title{\vspace{-1em}
Coherence Capacity in Modal Triplet Theory:\\[0.25em]
\large Normalized Margins, Metric Clearance, and the Invariance Class}
\author{Peter Nero}
\date{Version 4, July 2026}

\begin{document}
\maketitle

\begin{abstract}
Modal Triplet Theory (MTT) uses several independent control conditions:
spectral separation, projector and resolvent bounds, complementary
stability, leakage control, contraction, truncation accuracy, descent, and
domain or hyperbolicity requirements.  A quantitative compression of these
conditions must preserve their units, logical independence, and source
provenance.  A common zero set alone is insufficient because an arbitrary
reparameterization can change gradients, rates, magnitudes, and any proposed
physical coupling.

This paper defines a \emph{coherence-capacity record}.  Each declared
admissibility condition has a signed slack, a provenance, a tolerance, and
a positive scale with the same units.  Division by those scales produces a
dimensionless reserve vector.  Its minimum is the bottleneck capacity, while
a separately declared metric gives a geometric clearance from the
inadmissible set.  We prove four elementary but useful facts: positivity of
the bottleneck is exactly simultaneous satisfaction of the declared
conditions; its value gives a quantitative perturbation certificate;
uniform positive capacity is equivalent to a uniform reserve on a region;
and controlled changes of margin coordinates or bi-Lipschitz metrics
preserve a capacity \emph{class} up to explicit comparison constants.
Distance and slack definitions agree quantitatively only when an error-bound
or transversality estimate is supplied.

The scalar is therefore a normalized diagnostic compression, not a new
field, conserved charge, probability, force, entropy, or gravitational
coupling.  Capacity exhaustion means that at least one declared certificate
has reached its boundary.  It does not by itself imply failure of the upper
dynamics, nonexistence of a measurable section, irreversible evolution,
collapse, a horizon, or an arrow of time.  Those conclusions require
separate dynamics, constitutive laws, instruments, or geometric theorems.
The MTT Foundation remains the owner of the complete
admissibility ledger and continuation alternatives; this paper supplies the
metric and normalization layer needed to compare and transport that ledger
without conflating quantities of different units.
\end{abstract}

\section*{Revision note: Version 4}
\addcontentsline{toc}{section}{Revision note: Version 4}
\begin{description}[leftmargin=2.5cm,style=nextline]
\item[Supersedes]
Version 4 supersedes Version 3.0, DOI
\href{https://doi.org/10.5281/zenodo.18322057}
{10.5281/zenodo.18322057}.

\item[Reason]
Version 3 allowed any positive scalar with the desired zero set, asserted
that all MTT controls fail simultaneously, and inferred from scalar
exhaustion the absence of a global section, structural irreversibility,
Einstein gravity, an arrow of time, entropy, measurement collapse, and
undecidability.  A common zero set cannot support those conclusions.
Moreover, noninjectivity does not rule out a right inverse: an orthogonal
projection has the inclusion of its range as a bounded section.

\item[Resolution]
The primary object is now a sourced, dimensionless reserve vector.
The scalar capacity is its declared bottleneck compression, or a normalized
metric distance to the inadmissible set.  The metric, norms, tolerances,
scales, domain, and provenance are part of the definition.  ``Invariant''
now means equivalence up to explicit positive comparison constants under a
controlled change of record, not numerical uniqueness under arbitrary
reparameterization.

\item[Retained content]
The paper retains the useful idea that many independent MTT control
conditions can be summarized by their weakest remaining reserve.  It also
retains capacity exhaustion as a precise diagnostic first-exit event.

\item[Remaining boundary]
No universal choice of metric, normalization, or margin list is selected
for every MTT realization.  No transport, reset, probability, force,
thermodynamic, gravitational, or cosmological law is derived here.
\end{description}

\tableofcontents

\section{Introduction: the question in its correct form}
\label{sec:question}

An effective construction is rarely valid everywhere.  A spectral cluster
may cease to be isolated, a complementary resolvent may grow too large, a
contraction constant may reach one, or a truncation remainder may exceed its
tolerance.  These are different failures measured in different units.
Calling all of them ``loss of closure'' can be helpful prose, but it is not
yet a quantitative definition.

The corrected MTT Foundation records such conditions as separate margins
\cite{MTTFoundation2026}.  It also emphasizes that the entries can fail
independently and that a first exit from an admissible chart does not choose
what happens next.  Fixed Points I supplies the canonical conditional
spectral-projector and fixed-point framework on which several of those
margins are based \cite{FixedPointsI2026}.  This paper asks the narrower
question:
\begin{quote}
How can heterogeneous admissibility margins be compressed into one useful
number without erasing their units, provenance, or logical independence?
\end{quote}

There are two sensible answers.
\begin{enumerate}[leftmargin=2em]
\item Normalize every signed slack and retain the resulting vector.  Its
      least component is the active bottleneck.
\item Equip the state-and-control space with a declared metric and measure
      distance to the inadmissible set.
\end{enumerate}
Neither answer is canonical until its scales or metric are supplied.
Their numerical agreement is a theorem obligation, not a definition.

\subsection{What this paper does not attempt}

The present construction is a control diagnostic.  It does not derive the
underlying operator, action, state, physical geometry, or continuation law.
It does not turn a list of hypotheses into a new physical substance.  In
particular, a small reserve says that a declared approximation is fragile;
it does not say that nature exerts a force in the direction of increasing
reserve.

This distinction is familiar outside MTT.  Spectral perturbation theory
relates gaps to stability of invariant subspaces
\cite{Kato1966,DavisKahan1970}; Feshbach--Schur analysis relates block
operators and complementary resolvents to controlled effective operators
\cite{BachEtAl2003}.  A condition or error margin organizes those estimates.
It does not replace the operator theorem that produced them.

\section{The admissibility record}
\label{sec:record}

The capacity definition begins with a typed record, not with a scalar.

\begin{definition}[Admissibility-margin record]
\label{def:record}
An admissibility-margin record on a region \(U\) is
\[
 \cR_{\rm adm}
 =
 \bigl(
 U,d,\cI,\{s_i,\sigma_i,\cD_i,\pi_i\}_{i\in\cI}
 \bigr),
\]
where:
\begin{enumerate}[leftmargin=2.2em]
\item \(U\) is the state-and-control region on which all declared
      quantities can be compared.
\item \(d\) is a specified metric or extended metric on \(U\).  In an
      operator problem this may be an operator norm, graph norm, form
      metric, or a product metric, but the choice must be stated.
\item \(\cI\) is a finite set of required control rows.
\item \(s_i:\cD_i\to\mathbb R\) is a continuous signed slack on its declared
      domain \(\cD_i\subseteq U\).  Positive means that row \(i\) passes,
      zero is its boundary, and negative means that it fails.
\item \(\sigma_i>0\) is a declared reference scale with the same units as
      \(s_i\).  It is fixed independently of the point being certified, or
      its allowed state dependence and bounds are stated explicitly.
\item \(\pi_i\) is the provenance record for \(s_i\), its tolerance, norm,
      source theorem or assumption, uncertainty, and evaluation procedure.
\end{enumerate}
The common record domain is
\(\cD=\bigcap_{i\in\cI}\cD_i\), and the declared admissible set is
\[
 \cA
 =
 \{x\in\cD:s_i(x)>0\text{ for every }i\in\cI\}.
\]
\end{definition}

The strict inequality is useful when capacity is intended to measure
interior reserve.  A particular theorem may use nonnegative inequalities
and include part of the boundary; that convention must be recorded rather
than hidden in notation.

\subsection{Typical rows}

The complete Foundation ledger is longer than the following table, and it
remains the canonical source.  The table only shows how representative
conditions become signed slacks.

\begin{center}
\begin{tabularx}{\textwidth}{@{}L{0.20\textwidth}Y Y@{}}
\toprule
Row & Supplied control statement & Example signed slack \\
\midrule
Internal gap
& A selected spectral cluster has separation
  \(\gamma(x)>\gamma_{\min}\).
& \(s_{\rm gap}=\gamma-\gamma_{\min}\).\\

Resolvent or projector
& A declared contour or complementary resolvent obeys
  \(K(x)<K_{\max}\).
& \(s_{\rm res}=K_{\max}-K\).\\

Complementary stability
& The complementary semigroup has residual decay
  \(\omega_Q(x)>0\).
& \(s_Q=\omega_Q\).\\

Leakage or descent
& A leakage or commutator defect is below its certified tolerance:
  \(\ell(x)<\ell_{\max}\).
& \(s_{\rm leak}=\ell_{\max}-\ell\).\\

Contraction
& A map has a proved contraction factor \(q(x)<1\) on a declared invariant
  set and norm.
& \(s_{\rm ctr}=1-q\).\\

Truncation
& The complete remainder bound obeys
  \(E(x)<E_{\max}\).
& \(s_{\rm tr}=E_{\max}-E\).\\

Hyperbolicity or positivity
& A selected principal-symbol, Hessian, or Gram eigenvalue satisfies
  \(h(x)>h_{\min}\).
& \(s_{\rm hyp}=h-h_{\min}\).\\

Numerical certificate
& An interval inclusion, majorant, or residual test has a strict
  certified reserve \(\rho(x)>0\).
& \(s_{\rm cert}=\rho\).\\
\bottomrule
\end{tabularx}
\end{center}

The rows are not interchangeable.  A spectral gap does not prove
invariance; a contraction factor says nothing before its norm and invariant
set are fixed; and a small numerical residual is not an existence theorem
without the corresponding certification inequality.

\subsection{Why normalization is indispensable}

Suppose a gap is measured in inverse-length squared, an operator defect in
operator norm, and a truncation error in a dimensionless observable norm.
Taking their raw minimum is meaningless.  The positive scales
\(\sigma_i\) turn each slack into a dimensionless reserve
\[
 r_i(x)=\frac{s_i(x)}{\sigma_i}.
\]
The vector
\[
 r(x)=\bigl(r_i(x)\bigr)_{i\in\cI}
\]
is the primary capacity object.  It preserves which condition is limiting
and how each entry was normalized.

Scales can come from a theorem tolerance, a certified reference gap, an
instrument resolution, or a predeclared design budget.  Choosing a scale
after seeing which candidate one wishes to accept is target-dependent
selection, not an invariant definition.

\section{Bottleneck coherence capacity}
\label{sec:bottleneck}

\begin{definition}[Signed bottleneck and coherence capacity]
\label{def:capacity}
For a record as in Definition~\ref{def:record}, define
\[
 D_{\cR}(x)=\min_{i\in\cI}r_i(x),
 \qquad
 C_{\cR}(x)=\min\{1,\max\{0,D_{\cR}(x)\}\}.
\]
The active set is
\[
 I_{\rm act}(x)
 =
 \{i\in\cI:r_i(x)=D_{\cR}(x)\}.
\]
\end{definition}

The signed quantity \(D_{\cR}\) distinguishes interior, boundary, and
failure.  The clipped quantity \(C_{\cR}\in[0,1]\) is convenient for
comparison and visualization.  Clipping at one says only that reserves
larger than one declared scale are not distinguished by this summary.
The full vector should be retained in every archival calculation.

\begin{proposition}[Exact logical content of the bottleneck]
\label{prop:logical}
For every \(x\in\cD\):
\[
 x\in\cA
 \quad\Longleftrightarrow\quad
 D_{\cR}(x)>0
 \quad\Longleftrightarrow\quad
 C_{\cR}(x)>0.
\]
Moreover, \(D_{\cR}(x)=0\) means that at least one declared row is on its
boundary and none is negative.  It does not imply that every row vanishes.
\end{proposition}

\begin{proof}
The minimum of finitely many real numbers is positive exactly when every
number is positive.  Since every \(\sigma_i\) is positive,
\(r_i\) and \(s_i\) have the same sign.  The statement about the boundary
follows from the definition of a finite minimum.
\end{proof}

This small proposition corrects an important overstatement in Version 3.
The point of a bottleneck is precisely that one control may be exhausted
while the others retain large reserves.

\begin{proposition}[Perturbation certificate]
\label{prop:perturb}
Let \(x\in\cA\), and suppose \(y\in\cD\) satisfies
\[
 \|r(y)-r(x)\|_{\infty}<D_{\cR}(x).
\]
Then \(y\in\cA\).  More concretely, if every \(r_i\) is Lipschitz on a
neighborhood of \(x\) with constant \(L_i\), and
\(L=\max_iL_i\), then
\[
 d(x,y)<\frac{D_{\cR}(x)}{L}
 \quad\Longrightarrow\quad
 y\in\cA
\]
whenever \(L>0\) and the connecting neighborhood remains in \(\cD\).
\end{proposition}

\begin{proof}
For each \(i\),
\[
 r_i(y)
 >
 r_i(x)-D_{\cR}(x)
 \ge 0.
\]
The strict norm inequality makes every final inequality strict.  Under the
Lipschitz assumption,
\(\|r(y)-r(x)\|_\infty\le Ld(x,y)\), so the second statement follows.
\end{proof}

The proposition is the operational meaning of capacity: it is a certified
reserve in the chosen normalized coordinates.  It is not a statement about
energy or available work.

\begin{definition}[Regional capacity]
For \(K\subseteq\cD\), define
\[
 C_{\cR}(K)=\inf_{x\in K}C_{\cR}(x),
 \qquad
 D_{\cR}(K)=\inf_{x\in K}D_{\cR}(x).
\]
\end{definition}

\begin{proposition}[Uniform reserve]
\label{prop:uniform}
If \(D_{\cR}(K)>0\), then every declared margin has a uniform positive
normalized reserve on \(K\).  Conversely, if each row has a uniform lower
bound \(r_i(x)\ge\varepsilon_i>0\) on \(K\), then
\[
 D_{\cR}(K)\ge\min_i\varepsilon_i>0.
\]
\end{proposition}

\begin{proof}
The first implication follows from
\(r_i(x)\ge D_{\cR}(x)\ge D_{\cR}(K)\).  The converse follows by taking the
minimum first over \(i\) and then over \(x\in K\).
\end{proof}

Regional capacity is the correct quantity when a proof requires a
\emph{uniform} gap, resolvent, or contraction estimate.  A positive value at
one state does not certify an entire trajectory, slab, or parameter family.

\section{Metric clearance and distance to failure}
\label{sec:distance}

The bottleneck construction uses chosen constraint coordinates.  A
coordinate-free-looking alternative is distance to the complement, but it
still depends on a metric.

\begin{definition}[Metric clearance]
\label{def:clearance}
Assume \(\cA\) is open in the metric space \((\cD,d)\), and choose a positive
reference length \(L_d\) in the units of \(d\).  For \(x\in\cD\), define
\[
 \delta_d(x)=\dist_d(x,\cD\setminus\cA),
 \qquad
 C_d(x)=\min\left\{1,\frac{\delta_d(x)}{L_d}\right\}.
\]
\end{definition}

If the complement is empty, one may set \(\delta_d=+\infty\), but that
degenerate case contains no admissibility boundary and should be stated
separately.

\begin{proposition}[Basic metric properties]
\label{prop:distance}
The clearance \(\delta_d\) is \(1\)-Lipschitz:
\[
 |\delta_d(x)-\delta_d(y)|\le d(x,y).
\]
It is positive on \(\cA\) and zero on
\(\overline{\cD\setminus\cA}\).  If two metrics satisfy
\[
 a\,d(x,y)\le d'(x,y)\le b\,d(x,y)
\]
on the region of interest for constants \(0<a\le b<\infty\), then their
clearances satisfy
\[
 a\,\delta_d(x)\le\delta_{d'}(x)\le b\,\delta_d(x).
\]
\end{proposition}

\begin{proof}
For every \(z\in\cD\setminus\cA\), the triangle inequality gives
\(d(x,z)\le d(x,y)+d(y,z)\).  Taking infima over \(z\) yields
\(\delta_d(x)\le d(x,y)+\delta_d(y)\); exchanging \(x\) and \(y\) proves the
Lipschitz bound.  Openness gives positive distance locally only when the
point has a nonzero metric ball contained in \(\cA\), which is exactly the
definition.  The comparison statement follows by taking infima in the two
metric inequalities.
\end{proof}

\begin{remark}[Infinite-dimensional caution]
An open set guarantees that each interior point contains some metric ball,
so its pointwise clearance is positive.  It does not guarantee a positive
uniform clearance on a noncompact set.  Nor are an operator norm, graph
norm, strong topology, and weak topology interchangeable.  The topology and
metric used by the source theorem must be the ones used by the capacity
record.
\end{remark}

\subsection{When do slack and distance capacities agree?}

The two definitions have the same positive set by construction, but they
need not have comparable numerical values.  A flat margin such as
\(s(x)=\delta_d(x)^3\) and a steep margin such as
\(s(x)=\sqrt{\delta_d(x)}\) share a zero set while encoding very different
rates.

\begin{assumption}[Local error bound]
\label{ass:errorbound}
On a region \(V\subseteq\cA\), suppose there are constants
\(0<\alpha\le\beta<\infty\) such that
\[
 \alpha\,\frac{\delta_d(x)}{L_d}
 \le
 D_{\cR}(x)
 \le
 \beta\,\frac{\delta_d(x)}{L_d}
 \qquad (x\in V).
\]
\end{assumption}

\begin{proposition}[Slack--distance comparability]
\label{prop:comparison}
Under Assumption~\ref{ass:errorbound}, the clipped capacities obey
\[
 \alpha C_d(x)
 \le C_{\cR}(x)
 \le \beta C_d(x)
\]
on every subregion where both
\(\delta_d(x)/L_d\le1\) and \(D_{\cR}(x)\le1\).
More generally, the same inequalities hold locally after keeping the
unclipped signed bottleneck.
\end{proposition}

\begin{proof}
On the stated subregion no clipping occurs, so the claim is exactly
Assumption~\ref{ass:errorbound}.  The unclipped statement is the same comparison
without the subregion restriction.
\end{proof}

The hypothesis is substantial.  For affine inequalities, classical error
bounds give precisely this kind of comparison under suitable feasibility
conditions \cite{Hoffman1952}.  For nonlinear operator constraints one
needs a model-specific regularity, transversality, or metric-subregularity
estimate.  Sharing a boundary is not enough.

\section{What ``invariant'' can honestly mean}
\label{sec:invariance}

The word \emph{invariant} is useful only after the allowed transformations
are stated.

\begin{definition}[Capacity-equivalent records]
\label{def:equiv}
Two records \(\cR\) and \(\cR'\) on a region \(V\) are
capacity-equivalent if they define the same admissible set on \(V\) and
there exist constants \(0<a\le b<\infty\) such that
\[
 aD_{\cR}(x)\le D_{\cR'}(x)\le bD_{\cR}(x)
 \qquad
 \text{for every }x\in V\cap\cA.
\]
\end{definition}

This relation preserves positivity, exhaustion, and uniform-reserve
statements.  It does not assert that the numerical values are identical.

\begin{theorem}[Controlled reparameterization]
\label{thm:reparam}
Let \(r:V\to(0,\infty)^n\) be one normalized reserve vector and let
\(r'=F\circ r\) be another.  Suppose \(F\) preserves the positive cone and
there are constants \(0<a\le b<\infty\) such that
\[
 a\min_i u_i
 \le
 \min_jF_j(u)
 \le
 b\min_i u_i
\]
for every \(u\in r(V\cap\cA)\).  Then the two bottleneck records are
capacity-equivalent on \(V\).
\end{theorem}

\begin{proof}
Apply the displayed comparison to \(u=r(x)\).  The two minima are exactly
the two signed bottleneck capacities on the admissible region.
\end{proof}

Positive diagonal changes of normalization with scales bounded above and
below satisfy the theorem.  So do uniformly bi-Lipschitz changes of margin
coordinates that preserve the positive cone.  An arbitrary function
\(C'=f(C)\) need not.  For example, \(f(C)=C^3\) preserves the zero set but
has no positive linear lower comparison near zero.  It changes boundary
rates and makes \(\nabla C'\) vanish more rapidly.

\begin{remark}[The actual invariant]
The robust object is the equivalence class
\[
 [\cR]_{\rm cap}
\]
of records under controlled comparisons.  The positivity domain and
exhaustion boundary are qualitative invariants.  A numerical scalar is a
representative chosen by a metric and normalization convention.
\end{remark}

\section{Operator-theoretic examples}
\label{sec:examples}

The definitions become useful when the slacks come from real estimates.
The examples below illustrate the bookkeeping; they do not claim that one
universal list applies to every MTT model.

\subsection{A spectral projector}

Let \(A(x)\) be a self-adjoint operator family with a selected spectral
cluster separated from the rest of the spectrum by \(\gamma(x)\).
Let \(P(x)\) be the corresponding Riesz projector.  Perturbation theory
relates a nonzero gap and resolvent control to stability of the selected
subspace \cite{Kato1966,DavisKahan1970}.

A legitimate record could contain
\[
 s_{\rm gap}(x)=\gamma(x)-\gamma_{\min},
\]
and, for a declared contour \(\Gamma\),
\[
 K_\Gamma(x)
 =
 \sup_{z\in\Gamma}\|(z-A(x))^{-1}\|,
 \qquad
 s_{\rm res}(x)=K_{\max}-K_\Gamma(x).
\]
The record must state the operator domain, perturbation topology, contour,
target cluster, thresholds, and Sobolev or graph norm in which \(P(x)\) is
controlled.  The scalar capacity then reports which reserve is smaller.
It does not prove the Riesz formula or projector bound; those remain the
source theorem.

\subsection{A Schur--Feshbach truncation}

Let \(\cH=P\cH\oplus Q\cH\), \(Q=\Id-P\), and write a closed operator \(T\)
in blocks.  At a spectral parameter \(z\) for which \(QTQ-z\) is invertible
on its declared domain, the complementary correction has the form
\[
 R_F(z)
 =
 PTQ\,(QTQ-z)^{-1}QTP.
\]
If the effective construction requires
\(\|R_F(z)\|<\eta_{\max}\), the normalized truncation reserve may be
\[
 r_F(z)
 =
 1-
 \frac{
 \|PTQ\|\,
 \|(QTQ-z)^{-1}\|\,
 \|QTP\|
 }{\eta_{\max}}.
\]
This uses the actual three-factor control product.  The slogan
``mixing squared divided by a gap'' is only a special symmetric estimate
and must not replace the domain, resolvent, and block norms
\cite{BachEtAl2003}.

\subsection{A contraction argument}

Suppose \(F_x:K_x\to K_x\) is a contraction in a declared norm with factor
\(q(x)<1\).  One row may be
\[
 r_{\rm ctr}(x)=1-q(x).
\]
Another row must record whether \(F_x\) actually maps the chosen set into
itself, for example through a separately proved invariance reserve.  A
positive \(r_{\rm ctr}\) without a complete domain and invariant-set
statement is not a fixed-point certificate.  Conversely, once those
hypotheses are supplied, the Banach fixed-point theorem, not the word
``capacity'', proves existence and uniqueness.  Fixed Points I owns the
corresponding MTT framework \cite{FixedPointsI2026}.

\subsection{A numerical toy record}

Consider four already normalized rows at one point:
\[
 r(x)=(0.62,\;0.31,\;0.08,\;0.44).
\]
Then
\[
 D_{\cR}(x)=C_{\cR}(x)=0.08,
\qquad
 I_{\rm act}(x)=\{3\}.
\]
The record is admissible, but its third certificate is fragile.  Any
perturbation changing every normalized row by less than \(0.08\) preserves
admissibility by Proposition~\ref{prop:perturb}.  Nothing in these four numbers says
whether row 3 is a spectral gap, a truncation error, or a detector
threshold.  That meaning lives in the provenance \(\pi_3\).

\section{Capacity along an evolution}
\label{sec:dynamics}

Capacity can be evaluated along a trajectory without becoming the generator
of that trajectory.  Let \(x:[0,T]\to\cD\) be supplied by an upper or
effective evolution, and suppose the normalized rows
\(r_i(x(t))\) are absolutely continuous.

\begin{proposition}[First-exit reserve bound]
\label{prop:path}
If
\[
 \left|\frac{d}{dt}r_i(x(t))\right|\le v_i(t)
\]
for almost every \(t\), and \(v(t)=\max_i v_i(t)\), then
\[
 D_{\cR}(x(t))
 \ge
 D_{\cR}(x(0))-\int_0^t v(s)\,ds.
\]
Consequently, no declared margin can be exhausted before the integral of
the worst normalized rate reaches the initial bottleneck reserve.
\end{proposition}

\begin{proof}
For every \(i\),
\[
 r_i(x(t))
 \ge
 r_i(x(0))-\int_0^t v_i(s)\,ds
 \ge
 D_{\cR}(x(0))-\int_0^t v(s)\,ds.
\]
Taking the minimum over \(i\) proves the result.
\end{proof}

This is a useful monitoring bound.  It is not a conservation law.  A law
such as
\[
 \partial_t C+\nabla\cdot J_C=S_C
\]
requires a separately defined field, flux, source, regularity class, and
constitutive or variational principle.  Likewise, a barrier
\(-\log(\varepsilon+C)\) explicitly introduces a penalty into an algorithm;
its negative gradient is a chosen force in that model, not a consequence of
the capacity definition.

\subsection{What happens at zero?}

At a first time \(t_\ast\) with \(D_{\cR}(x(t_\ast))=0\), the record says
only which certificate has reached its declared boundary.  The corrected
Foundation distinguishes three possibilities:
\begin{enumerate}[leftmargin=2em]
\item stop the reduced description;
\item continue the upper dynamics and derive another reduced chart; or
\item add a reset map or reset kernel and thereby define a hybrid system.
\end{enumerate}
The boundary does not select among them.  If a reset has several outputs,
the boundary also does not provide their probabilities.

\section{Why the former physical deductions do not follow}
\label{sec:nodeductions}

The corrected definition is intentionally less dramatic than Version 3.
That is a gain in rigor: downstream physical claims can now be tested
against the exact additional data they require.

\subsection{A projection may have a section}

Let \(P:\cH\to P\cH\) be an orthogonal projection.  It is noninjective when
\(\ker P\ne\{0\}\), yet the inclusion
\[
 \iota:P\cH\hookrightarrow\cH
\]
is a bounded measurable map satisfying
\[
 P\circ\iota=\Id_{P\cH}.
\]
Thus noninjectivity, gap closure, or identification of upper states does
not by itself prove that no right inverse or measurable section exists.

A section selects one representative of each lower state.  It does not
recover the actual original upper state.  The impossible object for a
noninjective map is a \emph{left} inverse recovering every upper state:
\[
 L\circ P=\Id_{\cH}.
\]
The physically relevant questions are therefore whether a selected decoder
exists, whether it commutes with the dynamics on its domain, and whether it
recovers the information one claims.  Capacity exhaustion alone decides
none of these.

\subsection{Irreversibility and time}

The reduced chart can lose a certificate while the upper evolution remains
invertible.  It can also cross into another valid chart and later return.
Irreversibility requires a noninvertible lower semigroup, a coarse-graining
theorem, a reset rule, a record-stability mechanism, or a thermodynamic
limit.  A sequence of zero crossings is not, by itself, an arrow of time.

\subsection{Gravity}

One may postulate an effective action containing
\[
 \int\sqrt{-g}\,f(C)R,
\]
but this is extra dynamics.  Its variation generally contains derivatives
of \(f(C)\) and resembles a scalar-tensor coupling unless \(f(C)\) is fixed
and constant.  The capacity definition neither selects the action nor
derives the normalization \(1/(16\pi G)\).  Identifying
\(G^{-1}\) with an arbitrary representative of a capacity class is
especially ill-defined because \(C\mapsto C^3\) preserves the zero set while
changing the proposed coupling.

\subsection{Horizons, entropy, and measurement}

A horizon requires a Lorentzian geometry, causal structure, and field
equations.  An area law requires a state-counting, algebraic, or
entanglement argument and an independently derived coefficient.  A
measurement requires a system--apparatus interaction, pointer records, and
an outcome instrument.  Capacity can enter such models as one certified
control row, but it does not supply their missing structures.

\subsection{Probability and undecidability}

The location of a boundary does not define a probability measure over
possible continuations.  Nor does it imply undecidability.  An
undecidability theorem needs an explicit robust embedding of a universal
machine and a proof that its storage and operations remain admissible for
the required run.  Computational cost, unpredictability, first-exit
sensitivity, and formal undecidability are different claims.

\section{Relation to the MTT paper sequence}
\label{sec:sequence}

This revision gives the capacity sequence a coherent division of labor.

\begin{center}
\begin{tabularx}{\textwidth}{@{}L{0.28\textwidth}Y@{}}
\toprule
Document or topic & Correct role after this revision \\
\midrule
MTT Foundation
& Canonical owner of the complete admissibility ledger, independent
  hypotheses, and stop/continue/reset alternatives
  \cite{MTTFoundation2026}.\\

Fixed Points I--VI
& Canonical owners of their conditional projector, existence, stability,
  covariance, and synthesis results.  Capacity may summarize proved
  reserves but does not reproach or re-prove those theorems
  \cite{FixedPointsI2026,FixedPointsVI2026}.\\

This paper
& Owner of the normalized reserve record, bottleneck compression, metric
  clearance, comparison conditions, and capacity-equivalence language.\\

Controlled truncation
& Must emit the actual block-domain, resolvent, mixing, remainder, and error
  rows before capacity can summarize them.\\

Capacity dynamics
& May study a declared constitutive model, but conservation and transport
  do not follow from the margin definition.\\

Particles and forces
& May use capacity in an action or Hamiltonian if that coupling is sourced;
  \(\nabla C\) is not automatically a force.\\

Horizons and cosmology
& Require covariant dynamics, causal and state data, normalization, and
  comparison with observations.  Capacity language alone is interpretive.\\

Capacity-gated algorithms
& Are hybrid constrained systems.  Their flow domains, guards, reset maps
  or kernels, and invariants must be specified explicitly.\\
\bottomrule
\end{tabularx}
\end{center}

This order prevents circular reasoning.  A downstream paper may not choose
a scalar capacity to obtain a desired force, use that force to define the
dynamics, and then claim the resulting trajectory proves the original
capacity law.

\section{Source discipline and falsifiability}
\label{sec:source}

A capacity value is only as strong as its rows.  Every published value
should include:
\begin{enumerate}[leftmargin=2em]
\item the exact state, parameter point, and domain;
\item the finite list of required margins;
\item the norm, topology, tolerance, and units for each margin;
\item the normalization scale and why it was fixed;
\item the theorem, assumption, numerical certificate, or measured input
      sourcing each row;
\item interval or uncertainty bounds;
\item the active constraint set;
\item whether the claim is pointwise, local, regional, or trajectory-wide;
\item a hash-addressed calculation when numerical values are used; and
\item a continuation statement that is separate from the boundary test.
\end{enumerate}

The thresholds must be fixed before they are used to exclude data or select
a branch.  Otherwise the capacity is an after-the-fact score.  A held-out
failure of a claimed positive reserve falsifies that certificate or one of
its source assumptions.  It does not by itself falsify every MTT encoding.

\subsection{A minimal archival packet}

For machine-readable work, one capacity packet should contain at least the
fields
\begin{center}
\begin{tabularx}{0.94\textwidth}{@{}L{0.29\textwidth}Y@{}}
\texttt{record\_id}, \texttt{source\_commit}
& Identity and immutable source provenance.\\
\texttt{domain}, \texttt{metric}
& The common evaluation region and quantitative topology.\\
\texttt{rows}, \texttt{scales}, \texttt{intervals}
& Signed slacks, normalization choices, and certified uncertainty.\\
\texttt{active\_set}, \texttt{capacity}, \texttt{scope}
& The limiting rows, scalar compression, and claim domain.\\
\texttt{continuation\_policy}
& Stop, upper continuation, or a separately sourced reset rule.\\
\end{tabularx}
\end{center}
The scalar should never be stored without the row vector.  This makes a
future change of normalization auditable and prevents one paper's
diagnostic from silently becoming another paper's physical parameter.

\section{Limitations and open work}
\label{sec:limits}

The current construction closes a definitional problem, not a physical
source theorem.
\begin{enumerate}[leftmargin=2em]
\item There is no universal metric on the space of all MTT realizations.
\item Different operator problems naturally use different norms and
      domains.
\item Slack--distance equivalence needs a model-specific error bound.
\item A finite margin list is only complete relative to a declared theorem
      or calculation contract.
\item The capacity class does not select a preferred scalar representative.
\item No dynamics, probability, entropy, force, or geometry follows from
      the scalar without additional structure.
\end{enumerate}

These limits are productive.  They turn the word ``capacity'' into a
checklist of explicit obligations.  A future model can strengthen the
construction by selecting its metric and scales from the same upper source
as its operators, and then proving that the resulting record descends
naturally through the relevant projections.

\section{Conclusion}
\label{sec:conclusion}

Coherence capacity is best understood as the normalized clearance of a
declared effective construction from its nearest certified failure.
The full reserve vector is primary.  Its minimum is a useful bottleneck
summary, and metric distance is a useful geometric alternative.  Both are
quantitative only after their scales, metric, norms, and provenance are
fixed.

This corrected formulation preserves the original intuition while removing
claims the scalar could not support.  Capacity exhaustion identifies a
first failed certificate; it does not decide the fate of the upper state or
manufacture a new physical law.  The reward is a definition that can be
used consistently across spectral projectors, fixed-point arguments,
controlled truncations, numerical certificates, and future MTT
realizations.

\section*{Reproducibility statement}

This paper proves structural statements from the displayed definitions and
inequalities.  It reports no fitted parameter, numerical prediction, or
calculation-derived physical value.  The source TeX, bibliography, revision
audit, and generated PDF are versioned in the MTT papers repository.  Papers
that publish numerical capacity values must additionally provide the
machine-readable packet described in Section~\ref{sec:source}.

\bibliographystyle{plain}
\bibliography{main}

\end{document}
